市场对 H1 预告给出了两阶段反应：公告披露窗口内股价上行，随后快速回撤。这个过程说明持续性与现金质量已成为交易焦点，但不能被写成“公告直接导致涨停”的因果结论——CNINFO 的元数据确认 7 月 7 日为披露日，却未能确认精确盘中/盘后时间。

\section{事件窗口：高波动后的相对回撤扩大}

7 月 7 日雅化收盘上涨 10.0\%，其后至 7 月 23 日收于 17.75 元，较 7 月 7 日回撤 27.99\%。同期天齐、赣锋和盛新的平均回撤为 22.75\%，沪深 300 指数回撤 1.76\%。因此，雅化的波动既包含锂盐周期交易，也存在公司层面对 H2 持续性和利润质量的额外折价；但没有订单簿、资金席位或持仓的原始数据，本报告不对资金方向作结论。

\begin{exhibitbox}[表 7-1：业绩预告窗口的价格与成交位置]
\begin{tabularx}{\exhibitboxwidth}{L{3.5cm}R{3.0cm}X}
\toprule
指标 & 数值 & 解释边界 \\
\midrule
2026-07-23 收盘价 / 市值 & 17.75 元 / 204.58 亿元 & 总股本 11.526 亿股复算；作为估值分母 \\
7 月 7 日当日表现 & +10.0\% & 位于公告披露日；不宣称已证明盘中因果 \\
7 月 7 日至 7 月 23 日回撤 & -27.99\% & 价格路径事实，不等于资金流结论 \\
同行平均回撤 / 沪深 300 & -22.75\% / -1.76\% & 同行相对表现提示额外公司折价的可能性 \\
20 日价格区间 & 15.87--24.65 元 & 已发生波动区间，不是目标价或技术承诺 \\
\bottomrule
\end{tabularx}
\end{exhibitbox}
\sourcenote{东财行情原始快照与事件窗口日线；CNINFO 公告查询元数据；同行相对表现包。}

\section{估值拥挤度：低 PE 同时包含上行与现金折价}

以 17.75 元计算，三份完整券商报告的 2026E EPS 2.01--2.63 元对应 6.75--8.83 倍 PE；本报告基准 EPS 1.74 元对应约 10.20 倍。相对于三家锂企 12.79 倍的 2026E 平均 PE，雅化表面上存在折价，但该差异不能直接被解释为低估：同行的资源、产品和现金流结构不同，而雅化的经营现金流及 H1 单位经济尚待验证。

\begin{exhibitbox}[表 7-2：现价隐含的三种竞争性解释]
\begin{tabularx}{\exhibitboxwidth}{L{3.2cm}X X}
\toprule
解释 & 17.75 元可能反映什么 & 下一项证据 \\
\midrule
周期高点折价 & 市场不相信锂价、销量或价差能延续到 H2 & H1 收入、分部毛利、库存和后续销量 \\
现金质量折价 & 市场要求利润转为经营现金流与回款 & H1 OCF、应收、存货和项目结算说明 \\
持续性被低估 & 市场低估历史民爆底座与经营效率改善 & 民爆分部利润/现金与锂盐单位经济同时验证 \\
\bottomrule
\end{tabularx}
\end{exhibitbox}

市场锚在第八章占有限权重：它能反映当前交易可观察的定价，却不能取代自有基本面模型。若正式 H1 证实现金质量和单位经济，市场折价可收窄；若仅确认合并利润而现金继续弱，低倍数本身不构成加仓理由。
